The four fibre-optic probe locations are identified separately in each
ultrasound recording by image-based preprocessing. The first 20 frames are
averaged to reduce frame-to-frame variation. A vertical intensity projection
of this average frame is used to locate the brightest horizontal band; the
image is then cropped to a region extending 180 pixels above and below that
location (clipped at the image boundaries). Sensor candidates are extracted
from the averaged crop after contrast-limited adaptive histogram
equalization, Gaussian smoothing, and binary thresholding. Contours with
areas between 20 and 5000 pixels are retained, and candidates with an
enclosing-circle radius greater than 60 pixels are discarded.

If exactly four candidates remain, they are used directly. If more remain,
the four largest candidates among the 12 largest are selected as the group
with the highest axis-aligned rectangularity score. This score favours
similar opposing side lengths and diagonals, and horizontal and vertical
alignment. Fewer than four candidates do not yield a complete sensor
localisation. The selected points are ordered by image position: the upper
row from left to right is labelled M1 and M2, and the lower row from left to
right is labelled M3 and M4. Their center coordinates are stored in
\texttt{sensor\_coordinates.json} in pixels relative to the cropped video. A
visualization with the detected locations is also saved for inspection.

During dataset loading, these coordinates are scaled to the model input
resolution and used to place the corresponding M1--M4 temperature readings
in sparse target maps. Thus, the implemented procedure localises probes in
the cropped ultrasound image on a per-recording basis; it does not estimate
physical coordinates or a geometric transformation from the phantom to the
ultrasound frame.
